\documentclass[12pt,a4paper]{article}
\usepackage[utf8]{inputenc}
\usepackage[vietnamese]{babel}
\usepackage{geometry}

\usepackage{xcolor}
\usepackage{parskip}

% Khai báo các gói lệnh cần thiết
\usepackage[T5]{fontenc} % Font tiếng Việt
\usepackage{mathptmx} % Font Times New Roman giả lập
\usepackage{graphicx} % Chèn ảnh (logo)
\usepackage{tikz} % Vẽ khung bìa
\usetikzlibrary{calc}
\usepackage{anyfontsize} % Thay đổi cỡ chữ linh hoạt
\usepackage{hyperref} % Hỗ trợ chèn link web
\usepackage{indentfirst} % Ép thụt đầu dòng cho đoạn văn đầu tiên sau tiêu đề
\usepackage{minted}



% Cấu hình lề trang (Chuẩn A4)
\geometry{
  top=2cm,
  bottom=2cm,
  left=3cm,
  right=2cm
}

% Lệnh vẽ khung viền (2 nét: nét đậm ngoài, nét mảnh trong)
\newcommand{\drawCoverFrame}{
    \begin{tikzpicture}[remember picture, overlay]
        % Nét đậm bên ngoài
        \draw[line width=2pt] 
            ($(current page.north west) + (2.5cm,-2.0cm)$) 
            rectangle 
            ($(current page.south east) + (-1.5cm,2.0cm)$);
        % Nét mảnh bên trong
        \draw[line width=0.5pt] 
            ($(current page.north west) + (2.6cm,-2.1cm)$) 
            rectangle 
            ($(current page.south east) + (-1.6cm,2.1cm)$);
    \end{tikzpicture}
}

% =========================================================
% CẤU HÌNH ĐỊNH DẠNG TIÊU ĐỀ (Dùng gói titlesec)
% =========================================================
\usepackage{titlesec}

% Cho phép đánh số tiêu đề đến cấp 4 (để hiển thị 1.1.1.1)
\setcounter{secnumdepth}{4}
% Cho phép hiển thị cấp 4 vào Mục lục (nếu không cần hiện trong mục lục thì bỏ dòng này)
\setcounter{tocdepth}{4} 

% 1. Cấp Chương (\section): Sang trang, Căn giữa, In đậm. 
\titleformat{\section}[block]
  {\bfseries\Large}
  {Bài \thesection:}
  {0.5em}
  {}
% (Ghi chú: Bạn sẽ tự gõ IN HOA nội dung bên trong lệnh để tránh lỗi font tiếng Việt)


% 2. Cấp 1.1 (\subsection): In đậm
\titleformat{\subsection}
  {\bfseries\Large}
  {\thesubsection}
  {0.5em}
  {}
\titlespacing*{\subsection}{0.75cm}{1.5ex}{1ex} % Thụt lề 0.75cm

% 3. Cấp 1.1.1 (\subsubsection): In đậm, In nghiêng
\titleformat{\subsubsection}
  {\bfseries\itshape\Large}
  {\thesubsubsection}
  {0.5em}
  {}
\titlespacing*{\subsubsection}{1.5cm}{1.5ex}{1ex} % Thụt lề 1.5cm

% 4. Cấp 1.1.1.1 (\paragraph): In nghiêng
\titleformat{\paragraph}[hang]
  {\itshape\Large}
  {\theparagraph}
  {0.5em}
  {}
\titlespacing*{\paragraph}{2.25cm}{1.5ex}{1ex} % Thụt lề 2.25cm



\definecolor{backcolour}{rgb}{0.95,0.95,0.92}


\setminted{
    encoding=utf8,
    bgcolor=backcolour,
    fontsize=\footnotesize,
    breaklines=true,
    linenos=true,
    numbersep=5pt,
    tabsize=4
}

\begin{document}

% TRANG BÌA CHÍNH (Bìa 1 - CÓ KHUNG)
% ====================================================================
\begin{titlepage}
    \drawCoverFrame % <--- CÓ VẼ KHUNG
    \centering
    
    % Phần Header: Trường và Khoa
    \vspace*{0.5cm}
    {\fontsize{13}{15}\selectfont
    \textbf{ỦY BAN NHÂN DÂN THÀNH PHỐ HỒ CHÍ MINH}\\
    \textbf{TRƯỜNG ĐẠI HỌC SÀI GÒN}\\
    \textbf{KHOA TOÁN - ỨNG DỤNG}\\[0.5cm]
    -----oOo-----\\[0.75cm]
    }
    
    % Logo
    \includegraphics[width=4cm, keepaspectratio]{logoSGU.png} \\[1cm]
    
    % Phần Tên Báo Cáo
    {\fontsize{20}{24}\selectfont \textbf{BÀI TẬP}}\\[0.5cm]
    {\fontsize{20}{24}\selectfont \textbf{HỌC SÂU}}\\[0.5cm]
    
    % Tên Đề Tài
    {\fontsize{20}{24}\selectfont \textbf{ĐỀ TÀI: BÁO CÁO LAB01}}\\[1.25cm]
    
    % Thông tin Giảng viên và Sinh viên
    \begin{flushleft}
    \hspace{1cm} 
    \begin{minipage}{0.8\textwidth}
        \fontsize{16}{16}\selectfont
        \textbf{MÔN HỌC:} HỌC SÂU\\[0.75cm]
        \textbf{MÃ HỌC PHẦN:} 858048\\[0.75cm]
            \textbf{SINH VIÊN THỰC HIỆN:} TRƯƠNG TUYẾT NGHI \\[0.75cm]
        \textbf{LỚP:} KHOA HỌC DỮ LIỆU - DDU1231 \\[0.5cm]
        \textbf{GIẢNG VIÊN:} ĐỖ NHƯ TÀI\\[0.75cm]
    \end{minipage}
    \end{flushleft}
    
    \vfill
    
    % Địa điểm, Thời gian
    {\fontsize{13}{15}\selectfont
    {TP. Hồ Chí Minh, tháng 09, năm 2026}
    }
    \vspace*{1cm}
    
\end{titlepage}

% 
\newpage
\thispagestyle{empty}
% Đổi tên mặc định từ "Contents" thành "MỤC LỤC"
\renewcommand{\contentsname}{MỤC LỤC}
\begingroup
\titleformat{\section}[block]{\centering\bfseries\Large}{}{0pt}{}
\tableofcontents
\endgroup
\thispagestyle{empty}
\cleardoublepage


% BÀI 1
\section{Tạo Tensor cơ bản}
Tạo một Tensor số nguyên kích thước $2 \times 3$ và hiển thị kết quả.

\textbf{Code:}
\begin{minted}{python}
import torch

a = torch.tensor(
    [[1, 2, 3],
    [4, 5, 6]],
    dtype=torch.int32
)

print(a)
\end{minted}

\textbf{Output:}
\begin{minted}[fontfamily=tt, fontsize=\normalsize, bgcolor=none, linenos=false, formatcom=\color{gray}]{text}
tensor([[1, 2, 3],
        [4, 5, 6]], dtype=torch.int32)
\end{minted}

\textbf{Yêu cầu:}
Cho biết kích thước, số chiều và kiểu dữ liệu của Tensor a.


\textbf{Kết quả:}
\begin{itemize}
    \item Kích thước (Shape): $2 \times 3$.
    \item Số chiều (ndim): 2 chiều (gồm 2 hàng và 3 cột).
    \item Kiểu dữ liệu (dtype): Số nguyên (int32).
\end{itemize}


\section{Tạo các Tensor thường dùng}
Thực hành tạo Tensor ngẫu nhiên, Tensor toàn số 0, toàn số 1, ma trận đơn vị 

\textbf{Code:}
\begin{minted}{python}
import torch

a = torch.rand(3, 4)
b = torch.randn(3, 4)
c = torch.zeros(2, 3)
d = torch.ones(2, 3)
e = torch.eye(4)

print(a)
print(b)
print(c)
print(d)
print(e)
\end{minted}


\textbf{Output:}
\begin{minted}[fontfamily=tt, fontsize=\normalsize, bgcolor=none, linenos=false, formatcom=\color{gray}]{text}
tensor([[0.6564, 0.2357, 0.0230, 0.2447],
        [0.2288, 0.0760, 0.0233, 0.6651],
        [0.5633, 0.0643, 0.6013, 0.8193]])
tensor([[ 0.0746, -0.4413,  0.1134,  0.9300],
        [ 1.2199, -1.2352,  0.1010, -0.2774],
        [ 0.3536, -0.3021, -0.0517, -0.5805]])
tensor([[0., 0., 0.],
        [0., 0., 0.]])
tensor([[1., 1., 1.],
        [1., 1., 1.]])
tensor([[1., 0., 0., 0.],
        [0., 1., 0., 0.],
        [0., 0., 1., 0.],
        [0., 0., 0., 1.]])
\end{minted}

\textbf{Yêu cầu:}
\begin{itemize}
    \item Nhận xét sự khác nhau giữa torch.rand() và torch.randn().
    \item Tạo thêm một Tensor 5 × 5 chứa toàn giá trị 10.
\end{itemize}

\textbf{Kết quả:}
\begin{itemize}
    \item Sự khác nhau giữa torch.rand() và torch.randn():
    \begin{itemize}
        \item \texttt{torch.rand()}: Sử dụng phân phối đều $[0, 1)$, các giá trị trong ma trận chỉ nằm trong khoảng 0 đến 1, không chứa giá trị âm.
        \item \texttt{torch.randn()}: Sử dụng phân phối chuẩn (Gauss) $(-\infty, +\infty)$, giá trị tập trung quanh 0, càng xa 0 càng ít xuất hiện và có chứa giá trị âm.
    \end{itemize}
    \item Tensor 5 × 5 chứa toàn giá trị 10
\end{itemize}

\begin{minted}{python}
# Tạo thêm một Tensor 5 × 5 chứa toàn giá trị 10.
f = torch.full((5, 5), 10)

print(f)
\end{minted}


\begin{minted}[fontfamily=tt, fontsize=\normalsize, bgcolor=none, linenos=false, formatcom=\color{gray}]{text}
tensor([[10, 10, 10, 10, 10],
        [10, 10, 10, 10, 10],
        [10, 10, 10, 10, 10],
        [10, 10, 10, 10, 10],
        [10, 10, 10, 10, 10]])
\end{minted}





\section{Kiểm tra thông tin Tensor}
Sử dụng \texttt{shape}, \texttt{size()}, \texttt{ndim} và \texttt{dtype} để kiểm tra các thuộc tính của Tensor.

\textbf{Code:}
\begin{minted}{python}
import torch
a = torch.zeros(2, 3, 4)

print("Tensor:")
print(a)
print("Shape:", a.shape)
print("Size:", a.size())
print("Number of dimensions:", a.ndim)
\end{minted}

\textbf{Output:}
\begin{minted}[fontfamily=tt, fontsize=\normalsize, bgcolor=none, linenos=false, formatcom=\color{gray}]{text}
Tensor:
tensor([[[0., 0., 0., 0.],
         [0., 0., 0., 0.],
         [0., 0., 0., 0.]],

        [[0., 0., 0., 0.],
         [0., 0., 0., 0.],
         [0., 0., 0., 0.]]])
Shape: torch.Size([2, 3, 4])
Size: torch.Size([2, 3, 4])
Number of dimensions: 3
\end{minted}

\textbf{Yêu cầu:}
Giải thích ý nghĩa của \texttt{torch.Size([2, 3, 4])}.

\textbf{Kết quả:}
\begin{itemize}
    \item Tensor có kích thước $2 \times 3 \times 4$, gồm 3 chiều và tổng cộng $2 \times 3 \times 4 = 24$ phần tử.
    \item Tensor gồm 2 ma trận, mỗi ma trận có 3 hàng và 4 cột. Các phần tử đều bằng 0.
\end{itemize}

\section{Slicing và Indexing}
Truy xuất một phần dữ liệu từ Tensor nhiều chiều.

\textbf{Code:}
\begin{minted}{python}
import torch
a = torch.randn(3, 4, 5)
print("Original tensor:")
print(a)
print("\nTensor a[1]:")
print(a[1])
print("\nTensor a[1:, 2:4]:")
print(a[1:, 2:4])
\end{minted}

\textbf{Output:}
\begin{minted}[fontfamily=tt, fontsize=\normalsize, bgcolor=none, linenos=false, formatcom=\color{gray}]{text}
Original tensor:
tensor([[[ 0.1507, -1.8328, -1.6220, -0.2177,  0.9911],
         [ 1.0336,  0.3107,  0.1657,  0.9995, -0.2127],
         [ 0.6844,  0.9617,  1.3829,  0.7071, -0.0508],
         [-0.1017, -0.2587,  1.9031, -0.4535, -0.3880]],

        [[-0.9257,  0.2646, -0.7212, -0.5850,  0.8321],
         [-0.4802, -0.3190,  0.6407,  0.9775, -0.2529],
         [-0.1893, -0.3397, -0.0297, -1.0659, -0.9922],
         [ 0.6998, -0.1418, -0.3899, -0.1365,  3.0168]],

        [[ 0.3761, -0.9744,  0.5913,  1.0597, -0.0846],
         [ 1.3485,  2.4708,  0.6384, -0.6667,  1.5369],
         [ 0.4540, -0.0232,  0.3527,  0.1497, -0.3582],
         [-0.1223,  0.5231, -0.3172,  0.7384, -0.4826]]])

Tensor a[1]:
tensor([[-0.9257,  0.2646, -0.7212, -0.5850,  0.8321],
        [-0.4802, -0.3190,  0.6407,  0.9775, -0.2529],
        [-0.1893, -0.3397, -0.0297, -1.0659, -0.9922],
        [ 0.6998, -0.1418, -0.3899, -0.1365,  3.0168]])

Tensor a[1:, 2:4]:
tensor([[[-0.1893, -0.3397, -0.0297, -1.0659, -0.9922],
         [ 0.6998, -0.1418, -0.3899, -0.1365,  3.0168]],

        [[ 0.4540, -0.0232,  0.3527,  0.1497, -0.3582],
         [-0.1223,  0.5231, -0.3172,  0.7384, -0.4826]]])
\end{minted}

\textbf{Yêu cầu:}
\begin{itemize}
    \item Cho biết shape của \texttt{a[1]} và \texttt{a[1:, 2:4]}.
    \item Thử lấy hai phần tử cuối cùng trên chiều thứ ba.
\end{itemize}

\textbf{Kết quả:}
\begin{itemize}
    \item Với Tensor \texttt{a} có shape $(3,4,5)$, \texttt{a[1]} lấy ma trận thứ hai nên có shape $(4,5)$.

        \begin{minted}{python}
#Shape
print ("Shape a[1]: " ,(a[1].shape))

print ("Shape a[1:, 2:4]: " ,(a[1:, 2:4].shape))
    \end{minted}

    \begin{minted}[fontfamily=tt, fontsize=\normalsize, bgcolor=none, linenos=false, formatcom=\color{gray}]{text}
    Shape a[1]:  torch.Size([4])
    Shape a[1:, 2:4]:  torch.Size([2, 2])
    \end{minted}

    \item \texttt{a[1:, 2:4]} lấy hai ma trận cuối, giữ từ hàng 3  5 cột, nên có shape $(2,2,5)$.
    \begin{minted}{python}
# Lấy hai phần tử cuối cùng trên chiều thứ ba
a[:, :, -2:]
\end{minted}

    \begin{minted}[fontfamily=tt, fontsize=\normalsize, bgcolor=none, linenos=false, formatcom=\color{gray}]{text}
    tensor([[[-0.2177,  0.9911],
            [ 0.9995, -0.2127],
            [ 0.7071, -0.0508],
            [-0.4535, -0.3880]],

            [[-0.5850,  0.8321],
            [ 0.9775, -0.2529],
            [-1.0659, -0.9922],
            [-0.1365,  3.0168]],

            [[ 1.0597, -0.0846],
            [-0.6667,  1.5369],
            [ 0.1497, -0.3582],
            [ 0.7384, -0.4826]]])
    \end{minted}

\end{itemize}







\section{Reshape và Flatten}
Thay đổi hình dạng Tensor mà không làm thay đổi số phần tử.

\textbf{Code:}
\begin{minted}{python}
import torch
a = torch.randn(3, 4)
print("Original:")
print(a)
print("\nFlatten:")
print(a.ravel())
print("\nReshape to 3 x 2 x 2:")
print(a.reshape(3, 2, 2))
\end{minted}

\textbf{Output:}
\begin{minted}[fontfamily=tt, fontsize=\normalsize, bgcolor=none, linenos=false, formatcom=\color{gray}]{text}
Original:
tensor([[ 0.0294, -1.1352,  1.4799,  1.6398],
        [ 0.1498,  1.9367,  0.8845,  1.4105],
        [ 0.7305,  1.1645,  0.8977,  0.4912]])

Flatten:
tensor([ 0.0294, -1.1352,  1.4799,  1.6398,  0.1498,  1.9367,  0.8845,  1.4105,
         0.7305,  1.1645,  0.8977,  0.4912])

Reshape to 3 x 2 x 2:
tensor([[[ 0.0294, -1.1352],
         [ 1.4799,  1.6398]],

        [[ 0.1498,  1.9367],
         [ 0.8845,  1.4105]],

        [[ 0.7305,  1.1645],
         [ 0.8977,  0.4912]]])
\end{minted}

\textbf{Yêu cầu:}
\begin{itemize}
    \item Kiểm tra số phần tử trước và sau khi reshape.
    \item Thử reshape Tensor thành kích thước $2 \times 6$.
\end{itemize}

\textbf{Kết quả:}
\begin{itemize}
    \item Số phần tử trước và sau khi reshape.
    \begin{itemize}
        \item Số phần tử trước reshape: 12 phần tử, 
        \item Số phần tử sau reshape: 3 ma trận đều có 4 phần tử.
    \end{itemize}
    \begin{minted}{python}
# số phần tử trước và sau khi reshape.
print("Trước reshape:", a.shape)
print("Sau reshape:", a.reshape(3, 2, 2).shape)
\end{minted}

    \begin{minted}[fontfamily=tt, fontsize=\normalsize, bgcolor=none, linenos=false, formatcom=\color{gray}]{text}
    Trước reshape: torch.Size([3, 4])
    Sau reshape: torch.Size([3, 2, 2])
    \end{minted}
    
    \item Reshape Tensor thành kích thước $2 \times 6$.
    \begin{minted}{python}
# reshape Tensor thành kích thước 2 × 6
print(a.reshape( 2, 6))
\end{minted}

    \begin{minted}[fontfamily=tt, fontsize=\normalsize, bgcolor=none, linenos=false, formatcom=\color{gray}]{text}
    tensor([[ 0.0294, -1.1352,  1.4799,  1.6398,  0.1498,  1.9367],
            [ 0.8845,  1.4105,  0.7305,  1.1645,  0.8977,  0.4912]])
    \end{minted}
\end{itemize}







\section{Phép toán và nhân ma trận}
Thực hiện các phép toán phần tử và phép nhân ma trận bằng toán tử \texttt{@}.

\textbf{Code:}
\begin{minted}{python}
import torch

a = torch.randn(2, 3)
b = torch.randn(2, 3)

print("a:")
print(a)
print("b:")
print(b)

print("\na + b:")
print(a + b)
print("\na / b:")
print(a / b)

print("\na^2:")
print(a ** 2)

print("\nMatrix multiplication:")
print(a @ b.T)
\end{minted}

\textbf{Output:}
\begin{minted}[fontfamily=tt, fontsize=\normalsize, bgcolor=none, linenos=false, formatcom=\color{gray}]{text}
a:
tensor([[ 0.4173,  0.7319, -1.0494],
        [ 2.0330, -1.5014,  1.7079]])
b:
tensor([[ 0.4472,  0.6720,  0.9830],
        [-1.1775, -0.2348, -1.0270]])

a + b:
tensor([[ 0.8646,  1.4039, -0.0663],
        [ 0.8556, -1.7361,  0.6809]])

a / b:
tensor([[ 0.9331,  1.0891, -1.0674],
        [-1.7266,  6.3952, -1.6630]])

a^2:
tensor([[0.1742, 0.5357, 1.1011],
        [4.1332, 2.2541, 2.9168]])

Matrix multiplication:
tensor([[-0.3531,  0.4145],
        [ 1.5793, -3.7953]])
\end{minted}

\textbf{Yêu cầu:}
\begin{itemize}
    \item Phân biệt phép toán \texttt{a * b} với \texttt{a @ b.T}.
    \item Cho biết shape của kết quả \texttt{a @ b.T}.
\end{itemize}

\textbf{Kết quả:}
\begin{itemize}
    \item \texttt{a * b}: nhân từng cặp phần tử ở cùng vị trí. Hai Tensor có shape $(2,3)$ nên kết quả cũng có shape $(2,3)$.
    \item \texttt{a @ b.T}: là phép nhân ma trận. Tensor \texttt{b.T} có shape $(3,2)$ nên tích của ma trận $(2,3)$ với ma trận $(3,2)$ có shape $(2,2)$.

\begin{minted}{python}
# Ma trận a*b
print(a*b)

# Ma trận a @ b.T
print(a @ b.T)

#Shape a @ b.T
print("Shape a @ b.T: " ,(a @ b.T).shape)
    \end{minted}

    \begin{minted}[fontfamily=tt, fontsize=\normalsize, bgcolor=none, linenos=false, formatcom=\color{gray}]{text}
    tensor([[ 0.1867,  0.4918, -1.0316],
            [-2.3938,  0.3525, -1.7540]])
    tensor([[-0.3531,  0.4145],
            [ 1.5793, -3.7953]])
    Shape a @ b.T:  torch.Size([2, 2])
    \end{minted}
\end{itemize}


    

\section{Các hàm thống kê trên Tensor}
Tính trung bình, độ lệch chuẩn và tổng tích lũy theo từng cột.

\textbf{Code:}
\begin{minted}{python}
import torch
a = torch.randn(3, 4)
print(a)
print("\nMean:")
print(torch.mean(a, dim=0))
print("\nStandard deviation:")
print(torch.std(a, dim=0))
print("\nCumulative sum:")
print(torch.cumsum(a, dim=0))
\end{minted}

\textbf{Output:}
\begin{minted}[fontfamily=tt, fontsize=\normalsize, bgcolor=none, linenos=false, formatcom=\color{gray}]{text}
tensor([[ 0.2428, -2.0160, -1.1070, -0.3504],
        [-0.3763,  0.1236, -0.5444, -0.7849],
        [-0.4414, -1.3558,  0.2760,  1.2377]])

Mean:
tensor([-0.1916, -1.0827, -0.4585,  0.0341])

Standard deviation:
tensor([0.3777, 1.0956, 0.6955, 1.0647])

Cumulative sum:
tensor([[ 0.2428, -2.0160, -1.1070, -0.3504],
        [-0.1335, -1.8924, -1.6515, -1.1353],
        [-0.5749, -3.2482, -1.3755,  0.1023]])
\end{minted}

\textbf{Yêu cầu:}
Thay \texttt{dim=0} bằng \texttt{dim=1} và nhận xét sự khác nhau.

\textbf{Kết quả:}
\begin{itemize}
    \item \texttt{dim=0} tính theo từng cột, nên \texttt{mean} và \texttt{std} trả về 4 giá trị; \texttt{cumsum} cộng dồn từ trên xuống.
    \item \texttt{dim=1} tính theo từng hàng, nên \texttt{mean} và \texttt{std} trả về 3 giá trị; \texttt{cumsum} cộng dồn từ trái sang phải.
    \item Trong cả hai trường hợp, \texttt{cumsum} giữ nguyên shape $(3,4)$ của Tensor đầu vào.

    \begin{minted}{python}
# dim = 1
print(torch.mean(a, dim=1))

print("\nStandard deviation:")
print(torch.std(a, dim=1))

print("\nCumulative sum:")
print(torch.cumsum(a, dim=1))
\end{minted}

    \begin{minted}[fontfamily=tt, fontsize=\normalsize, bgcolor=none, linenos=false, formatcom=\color{gray}]{text}
    tensor([-0.8077, -0.3955, -0.0709])

    Standard deviation:
    tensor([0.9768, 0.3846, 1.0986])

    Cumulative sum:
    tensor([[ 0.2428, -1.7732, -2.8803, -3.2307],
            [-0.3763, -0.2527, -0.7972, -1.5821],
            [-0.4414, -1.7972, -1.5212, -0.2836]])
    \end{minted}
\end{itemize}



\section{Autograd cơ bản}
Tạo biến có \texttt{requires\_grad=True} và dùng \texttt{backward()} để tính đạo hàm tự động.

\textbf{Code:}
\begin{minted}{python}
import torch
x = torch.tensor(
3.6,
requires_grad=True
)
y = x * x
y.backward()
print("x =", x)
print("y =", y)
print("x.grad =", x.grad)
\end{minted}

\textbf{Output:}
\begin{minted}[fontfamily=tt, fontsize=\normalsize, bgcolor=none, linenos=false, formatcom=\color{gray}]{text}
x = tensor(3.6000, requires_grad=True)
y = tensor(12.9600, grad_fn=<MulBackward0>)
x.grad = tensor(7.2000)
\end{minted}

\textbf{Yêu cầu:}
\begin{itemize}
    \item Tính thủ công đạo hàm của $y=x^2$ tại $x=3.6$ và so sánh với \texttt{x.grad}.
    \item Đổi hàm thành $y=3x^2+2x+1$ và kiểm tra gradient.
\end{itemize}

\textbf{Kết quả:}
\begin{itemize}
    \item Với $y=x^2$, ta có $y'=2x$. Tại $x=3.6$, đạo hàm bằng $7.2$, khớp với giá trị \texttt{tensor(7.2000)} được in ra.
    \item Với $y=3x^2+2x+1$, ta có $y'=6x+2$. Tại $x=3.6$, đạo hàm bằng $23.6$. Với x.grad() = 23.6000, kết quả bằng nhau.

\begin{minted}{python}
# y=3x^2+2x+1, x= 3.6
x = torch.tensor(
3.6,
requires_grad=True
)
y = 3*x**2 + 2*x + 1
y.backward()
print("x =", x)
print("y =", y)
print("x.grad =", x.grad)
\end{minted}

    \begin{minted}[fontfamily=tt, fontsize=\normalsize, bgcolor=none, linenos=false, formatcom=\color{gray}]{text}
    x = tensor(3.6000, requires_grad=True)
    y = tensor(47.0800, grad_fn=<AddBackward0>)
    x.grad = tensor(23.6000)
    \end{minted}
\end{itemize}


\section{Gradient Descent tìm hệ số đa thức}
Sinh dữ liệu từ đa thức $y=x^2+2x+3$, sau đó dùng NAdam và autograd để tìm lại ba hệ số.

\textbf{Code:}
\begin{minted}{python}
import numpy as np
import torch

polynomial = np.poly1d([1, 2, 3])
N = 20

X = np.random.randn(N, 1) * 5
Y = polynomial(X)

XX = np.hstack([
X * X,
X,
np.ones_like(X)
])

w = torch.randn(3, 1, requires_grad=True)
x = torch.tensor(XX, dtype=torch.float32)
y = torch.tensor(Y, dtype=torch.float32)

optimizer = torch.optim.NAdam([w], lr=0.01)

print("Initial coefficients:")
print(w)

for epoch in range(1000):
    optimizer.zero_grad()
    y_pred = x @ w
    mse = torch.mean(torch.square(y - y_pred))
    mse.backward()
    optimizer.step()

print("Learned coefficients:")
print(w)
\end{minted}

\textbf{Output:}
\begin{minted}[fontfamily=tt, fontsize=\normalsize, bgcolor=none, linenos=false, formatcom=\color{gray}]{text}
Initial coefficients:
tensor([[-1.8516],
        [ 0.9187],
        [-0.0345]], requires_grad=True)
Learned coefficients:
tensor([[1.0014],
        [2.0022],
        [2.8323]], requires_grad=True)
\end{minted}

\textbf{Yêu cầu:}
\begin{itemize}
    \item So sánh ba hệ số học được với $[1,2,3]$.
    \item Giải thích \texttt{optimizer.zero\_grad()}, \texttt{mse.backward()} và \texttt{optimizer.step()}.
\end{itemize}

\textbf{Kết quả:}
\begin{itemize}
    \item So sánh ba hệ số học được với $[1,2,3]$.

    \begin{itemize}
        \item Theo output, ba hệ số học được là $[1.0014,\ 2.0022,\ 2.8323]$, tương ứng với các hệ số của $x^2$, $x$ và hệ số tự do.
        \item So với $[1,2,3]$, sai lệch xấp xỉ $0.0014$, $0.0022$ và $0.1677$. Hai hệ số đầu đã gần giá trị ban đầu hơn hệ số tự do; sau 1000 epoch, các hệ số chưa trùng hoàn toàn với đa thức ban đầu.
    \end{itemize}

    \item Giải thích \texttt{optimizer.zero\_grad()}, \texttt{mse.backward()} và \texttt{optimizer.step()}.
    \begin{itemize}
        \item \texttt{optimizer.zero\_grad()}: Reset toàn bộ đạo hàm của các biến hiện tại về ban đầu trước khi tính đạo hàm mới.
        \item \texttt{mse.backward()}: Tính đạo hàm của hàm mất mát (MSE) theo các hệ số w. 
        \item \texttt{optimizer.step()}: ập nhật hệ số w theo thuật toán NAdam để giảm sai số MSE.
\end{itemize}


\section{Giải hệ phương trình bằng Autograd}
Tối thiểu hóa tổng sai số bình phương để tìm $A,B,C,D$ thỏa mãn bốn phương trình

\textbf{Code:}
\begin{minted}{python}
import random
import torch

A = torch.tensor(random.random(), requires_grad=True)
B = torch.tensor(random.random(), requires_grad=True)
C = torch.tensor(random.random(), requires_grad=True)
D = torch.tensor(random.random(), requires_grad=True)

EPOCHS = 2000
optimizer = torch.optim.NAdam([A, B, C, D], lr=0.01)

for _ in range(EPOCHS):
    y1 = A + B - 9
    y2 = C - D - 1
    y3 = A + C - 8
    y4 = B - D - 2

    sqerr = y1*y1 + y2*y2 + y3*y3 + y4*y4

    optimizer.zero_grad()
    sqerr.backward()
    optimizer.step()

print("A =", A)
print("B =", B)

print("C =", C)
print("D =", D)
\end{minted}

\textbf{Output:}
\begin{minted}[fontfamily=tt, fontsize=\normalsize, bgcolor=none, linenos=false, formatcom=\color{gray}]{text}
A = tensor(4.6465, requires_grad=True)
B = tensor(4.3534, requires_grad=True)
C = tensor(3.3535, requires_grad=True)
D = tensor(2.3534, requires_grad=True)
\end{minted}

\textbf{Yêu cầu:}
\begin{itemize}
    \item Kiểm tra nghiệm thu được với bốn phương trình đã cho.
    \item Nhận xét vì sao bài toán có thể có nhiều nghiệm.
\end{itemize}

\textbf{Kết quả:}
\begin{itemize}
    \item Kiểm tra nghiệm thu được với bốn phương trình đã cho:
    
    $A+B=8.9999$, $C-D=1.0001$, $A+C=8.0000$, $B-D=2.0000$. 
    
    Sai số rất nhỏ.

    \item Bài toán có thể có nhiều nghiệm:

    Lấy phương trình thứ nhất cộng phương trình thứ hai rồi trừ phương trình thứ ba thu được $B-D=2$, chính là phương trình thứ tư. Do đó hệ chỉ có 3 phương trình độc lập với 4 ẩn.

    Đặt $D=t$, nghiệm tổng quát là $A=7-t$, $B=t+2$, $C=t+1$, $D=t$ với $t$ tùy ý thuộc tập số thực. Hệ có vô số nghiệm nên kết quả tối ưu là một nghiệm gần đúng trong họ nghiệm này.
\end{itemize}
    
    

\section{Bài tập tổng hợp}
\textbf{Yêu cầu:}
\begin{itemize}
    \item Tạo Tensor \texttt{X} có 24 phần tử và reshape lần lượt thành $4\times6$, $2\times12$ và $2\times3\times4$.
    \item Tạo hai Tensor \texttt{A}, \texttt{B} kích thước $3\times4$; tính \texttt{A+B}, \texttt{A*B}, \texttt{A/B}, trung bình theo cột và độ lệch chuẩn theo cột.
    \item Cho $y=2x^2+5x+3$. Dùng autograd tính $dy/dx$ tại $x=2$ và kiểm tra bằng phép tính thủ công.
    \item Thay đa thức trong Bài 9 bằng $y=2x^2-3x+5$; huấn luyện và cho biết các hệ số tìm được.
\end{itemize}

\textbf{Kết quả:}

\textbf{1. Tạo Tensor X có 24 phần tử và reshape lần lượt thành 4 × 6, 2 × 12 và 2 × 3 × 4.}
\begin{minted}{python}
# Tạo Tensor X có 24 phần tử và reshape lần lượt thành 4 × 6, 2 × 12 và 2 × 3 × 4.
X = torch.randn(24)
print("Tensor X:")
print(X)

print("Tensor 4x6")
print(X.reshape(4, 6))

print("Tensor 2x12")
print(X.reshape(2, 12))

print("Tensor 2x3x4")
print(X.reshape(2, 3, 4))
\end{minted}


\begin{minted}[fontfamily=tt, fontsize=\normalsize, bgcolor=none, linenos=false, formatcom=\color{gray}]{text}
Tensor X:
tensor([ 0.4969,  0.1460, -0.7639, -0.8069,  0.5555, -0.9160, -0.3864, -0.9593,
        -1.1049,  0.8254,  0.3225, -0.6399,  0.5124,  0.7560, -0.3735, -0.6824,
         0.4968,  1.2440, -0.1487, -0.6026,  0.9703,  0.3479,  2.0880,  1.0546])
Tensor 4x6
tensor([[ 0.4969,  0.1460, -0.7639, -0.8069,  0.5555, -0.9160],
        [-0.3864, -0.9593, -1.1049,  0.8254,  0.3225, -0.6399],
        [ 0.5124,  0.7560, -0.3735, -0.6824,  0.4968,  1.2440],
        [-0.1487, -0.6026,  0.9703,  0.3479,  2.0880,  1.0546]])
Tensor 2x12
tensor([[ 0.4969,  0.1460, -0.7639, -0.8069,  0.5555, -0.9160, -0.3864, -0.9593,
         -1.1049,  0.8254,  0.3225, -0.6399],
        [ 0.5124,  0.7560, -0.3735, -0.6824,  0.4968,  1.2440, -0.1487, -0.6026,
          0.9703,  0.3479,  2.0880,  1.0546]])
Tensor 2x3x4
tensor([[[ 0.4969,  0.1460, -0.7639, -0.8069],
         [ 0.5555, -0.9160, -0.3864, -0.9593],
         [-1.1049,  0.8254,  0.3225, -0.6399]],

        [[ 0.5124,  0.7560, -0.3735, -0.6824],
         [ 0.4968,  1.2440, -0.1487, -0.6026],
         [ 0.9703,  0.3479,  2.0880,  1.0546]]])
\end{minted}

Cả ba shape đều bảo toàn số phần tử: $4\times6=2\times12=2\times3\times4=24$.

\textbf{2. Tạo hai Tensor A, B kích thước 3 × 4; tính A+B, A*B, A/B, trung bình theo cột và độ lệch chuẩn theo cột.}

\begin{minted}{python}
# Tạo hai Tensor A, B kích thước 3 × 4; tính A+B, A*B, A/B, trung bình theo cột và độ lệch chuẩn theo cột.
A = torch.randn(3, 4)
B = torch.randn(3, 4)

print("A =", A)
print("B =", B)

print("\nA + B:", A + B)
print("A * B:", A * B)
print("A / B:", A / B)

print("Trung bình theo cột của A:", A.mean(dim=0))
print("Độ lệch chuẩn theo cột của A:", A.std(dim=0))
\end{minted}

\begin{minted}[fontfamily=tt, fontsize=\normalsize, bgcolor=none, linenos=false, formatcom=\color{gray}]{text}
A = tensor([[ 0.0352, -0.4872,  1.7132, -0.4583],
        [ 1.0278, -0.1278,  0.5428, -0.9981],
        [-0.0067,  1.0875, -0.2630,  0.1779]])
B = tensor([[ 0.5839,  0.4976,  1.9256, -0.5047],
        [ 0.4847,  0.5146,  1.9256, -0.2380],
        [-0.0819, -0.3730, -2.2113,  1.4797]])

A + B: tensor([[ 0.6191,  0.0104,  3.6388, -0.9630],
        [ 1.5125,  0.3868,  2.4683, -1.2361],
        [-0.0886,  0.7145, -2.4742,  1.6576]])
A * B: tensor([[ 2.0560e-02, -2.4244e-01,  3.2990e+00,  2.3129e-01],
        [ 4.9815e-01, -6.5752e-02,  1.0451e+00,  2.3752e-01],
        [ 5.4959e-04, -4.0567e-01,  5.8150e-01,  2.6322e-01]])
A / B: tensor([[ 0.0603, -0.9790,  0.8897,  0.9082],
        [ 2.1207, -0.2483,  0.2819,  4.1946],
        [ 0.0820, -2.9156,  0.1189,  0.1202]])
Trung bình theo cột của A: tensor([ 0.3521,  0.1575,  0.6643, -0.4262])
Độ lệch chuẩn theo cột của A: tensor([0.5856, 0.8252, 0.9937, 0.5887])
\end{minted}

\textbf{3. Cho y = 2x2 + 5x + 3. Dùng autograd tính dy/dx tại x = 2 và kiểm tra bằng phép tính thủ công.}

\begin{minted}{python}
#Cho y = 2x2 + 5x + 3. Dùng autograd tính dy/dx tại x = 2 và kiểm tra bằng phép tính thủ công.
x = torch.tensor(2.0, requires_grad=True)
y = 2 * x**2 + 5*x + 3
y.backward()
print("x=",x)
print("y=",y)
print("x.grad:", x.grad)
\end{minted}

\begin{minted}[fontfamily=tt, fontsize=\normalsize, bgcolor=none, linenos=false, formatcom=\color{gray}]{text}
x= tensor(2., requires_grad=True)
y= tensor(21., grad_fn=<AddBackward0>)
x.grad: tensor(13.)
\end{minted}

Với $y=2x^2+5x+3$, ta có $dy/dx=4x+5$. Tại $x=2$, $y=21$ và $dy/dx=13$, khớp với output \texttt{x.grad: tensor(13.)}.

\textbf{4. Huấn luyện đa thức $y=2x^2-3x+5$}

\begin{minted}{python}
# Thay đa thức trong Bài 9 bằng y = 2x2 - 3x + 5; huấn luyện và cho biết các hệ số tìm được.
import numpy as np
import torch

polynomial = np.poly1d([2, -3, 5])
N = 20

X = np.random.randn(N, 1) * 5
Y = polynomial(X)

XX = np.hstack([
X * X,
X,
np.ones_like(X)
])

w = torch.randn(3, 1, requires_grad=True)
x = torch.tensor(XX, dtype=torch.float32)
y = torch.tensor(Y, dtype=torch.float32)

optimizer = torch.optim.NAdam([w], lr=0.01)

print("Initial coefficients:")
print(w)

for epoch in range(1000):
    optimizer.zero_grad()
    y_pred = x @ w
    mse = torch.mean(torch.square(y - y_pred))
    mse.backward()
    optimizer.step()

print("Learned coefficients:")
print(w)
\end{minted}


\begin{minted}[fontfamily=tt, fontsize=\normalsize, bgcolor=none, linenos=false, formatcom=\color{gray}]{text}
Initial coefficients:
tensor([[-0.2890],
        [-0.0745],
        [ 1.3110]], requires_grad=True)
Learned coefficients:
tensor([[ 2.0145],
        [-3.0065],
        [ 3.9068]], requires_grad=True)
\end{minted}

\textbf{Nhận xét:}
\begin{itemize}
    \item Output là các hệ số $[2.0145,\ -3.0065,\ 3.9068]$. So với giá trị đích $[2,-3,5]$, sai lệch lần lượt xấp xỉ $0.0145$, $0.0065$ và $1.0932$.
    \item Hai hệ số của $x^2$ và $x$ đã khá gần giá trị đích, nhưng hệ số tự do còn lệch đáng kể so với 5. Có thể quá trình tối ưu chưa hội tụ sau 1000 epoch; sự khác biệt về độ lớn giữa các cột $X^2$, $X$ và cột toàn 1 cũng có thể ảnh hưởng đến tốc độ học từng hệ số.
\end{itemize}

\end{document}
